\section{Model Architecture Details}
\label{appendix:architecture}

\paragraph{\molm{}.}
\molm{} processes SMILES sequences through a series of blocks alternating between self-attention and feed-forward connections. The self-attention mechanism transforms input token sequences into query, key, and value representations, enabling the model to capture contextual relationships across the sequence. To encode positional information efficiently, \molm{} employs a modified RoFormer attention mechanism, which introduces position-dependent rotations $R_m$ of the queries and keys at each position $m$. The resulting attention at position $m$ is computed as:

\begin{equation}
\text{Attention}_m(Q, K, V) = \frac{\sum_{n=1}^{N} \langle \varphi(R_m q_m),\, \varphi(R_n k_n) \rangle\, v_n}{\sum_{n=1}^{N} \langle \varphi(R_m q_m),\, \varphi(R_n k_n) \rangle}
\end{equation}

where $\varphi$ is a random feature map that enables linear attention, avoiding the quadratic cost of standard softmax attention. Each token is embedded in a 768-dimensional space, and sequences are capped at 202 tokens, a threshold that accommodates over 99.4\% of molecules in the training corpus.

\paragraph{\smi{}.}
\smi{} employs the same linear RoFormer attention mechanism as \molm{} in its encoder, with a hidden size of 768, 12 attention heads, and 12 layers. The encoder-decoder architecture maps token embeddings into a compressed latent representation $\mathbf{z}$ via:

\begin{equation}
\mathbf{z} = \left(\text{LayerNorm}\left(\text{GELU}\left(\mathbf{x} W_1 + b_1\right)\right)\right) W_2
\end{equation}

and reconstructs the token sequence from $\mathbf{z}$ through a symmetric immersion operation followed by a language decoder layer that projects onto vocabulary logits for next-token prediction. The encoder contains 47 million parameters and the decoder 242 million parameters, totalling 289 million parameters in the base model.